\documentclass[a4paper,12pt]{scrartcl}
\usepackage{helvet}
\renewcommand{\familydefault}{\sfdefault}
\usepackage[a4paper,left = 2.5cm, right = 2.5cm,bottom=2.5cm, top=2.5cm]{geometry}
\usepackage[utf8]{inputenc}
\usepackage[T1]{fontenc}
\renewcommand{\thefigure}{\thesection.\arabic{figure}}
\setlength{\parindent}{0pt}
\usepackage{enumitem}
\usepackage{listings}
\usepackage{microtype} % hübschere Umbrüche
\lstset{
	basicstyle=\ttfamily\small,
	breaklines=true,
	breakatwhitespace=false,
	columns=fullflexible,
	keepspaces=true,
	upquote=true,
}
\newcommand{\code}[1]{\lstinline!#1!} % Kurzform für Inline-Code
\usepackage{csquotes}
\usepackage{graphicx}
\usepackage{wrapfig}
\usepackage{hyperref}
\usepackage{ifthen}
\usepackage{setspace}
\usepackage[backend=biber,style=numeric,sorting=nyt]{biblatex}
\DeclareNameAlias{default}{family-given}
\addbibresource{Literatur.bib}
\thispagestyle{plain}
\hypersetup{pageanchor=false}
\setcounter{tocdepth}{3}
\usepackage{float}
\usepackage{amsthm}
\usepackage{amsmath}
\newtheorem{definition}{Definition}[section]
\usepackage{ragged2e}
\usepackage{listings}
\usepackage{xcolor}
\lstset{
	basicstyle=\ttfamily\small,
	keywordstyle=\color{blue},
	commentstyle=\color{gray},
	stringstyle=\color{orange},
	numbers=left,
	numberstyle=\tiny,
	numbersep=5pt,
	frame=single,
	breaklines=true,
	captionpos=b
}
\usepackage[ngerman]{babel}
\usepackage[automark,headsepline,footsepline]{scrlayer-scrpage}
\clearpairofpagestyles
\automark[section]{section}
\ohead{}               % links leer
\ihead{\rightmark}  
\ofoot[\pagemark]{\pagemark}
\ifoot[Automatisiertes Modell- Hochregallager]{Automatisiertes Modell- Hochregallager}
\pagestyle{scrheadings}
\ModifyLayer[addvoffset=-1ex]{scrheadings.foot.above.line}
%\usepackage{chngcntr}
%\counterwithin{figure}{section}
%\counterwithin{table}{section}
%Kapitel + Abschnitt in Markierungen übernehmen
\begin{document}
	\pagenumbering{gobble}
	\begin{titlepage}
		\enlargethispage{4.0cm}
		\sffamily 								% Serifenlose Grundschrift für die Titelseite einstellen
		\parbox{0.5\linewidth}{
			\begin{flushleft}
				\includegraphics[width=0.4\linewidth]{images/DHBW_d_R_FN_46mm_4c.pdf}\\[5ex]
			\end{flushleft}
		}
		
		
		\begin{center}
			
			{\fontsize{20.74pt}{24pt}\selectfont
				\textbf{Automatisiertes Modell-Hochregallager}\\[1.5ex]}
			{\fontsize{14pt}{17pt}\selectfont
				\textbf{}\\[5ex]}
			{\fontsize{17pt}{20pt}\selectfont
				\textbf{T3\_3100}\\[2ex]}
			{\fontsize{14pt}{17pt}\selectfont
				Elektrotechnik\\[2ex]}
			{\fontsize{12pt}{14pt}\selectfont
				Automation\\[1ex]
				Duale Hochschule Baden-Württemberg Ravensburg, Campus Friedrichshafen\\[5ex]
				von\\[1ex]
				Maximilian Zorko und Willi Schaal\\[15ex]}
			
			
		\end{center}
		
		\begin{flushleft}
			{\fontsize{12pt}{14pt}\selectfont
				\begin{tabular}{ll}
					Abgabedatum:					& \quad  \\
					Bearbeitungszeitraum:		   	& \quad 16.10.2025 -    \\ 
					Matrikelnummer: 			& \quad 3960407/5732737 \\ 
					Kurs: 							& \quad TEA 23 \\
					% entfällt bei Studienarbeit
					Betreuerin / Betreuer:  & \quad Thorsten Kever \\ % Betreuerin / Betreuer der Arbeit
					%	Gutachterin / Gutachter: & \quad \gutachterdhbw \\ [2ex] % Gutachterin / Gutachter der DHBW (nur bei der Bachelorarbeit erforderlich)
				\end{tabular}
			}
		\end{flushleft}
		%%%%% Nachfolgende Zeilen einkommentieren, wenn Copyrightvermerk gewünscht ist
		\begin{flushleft}
			{\fontsize{11pt}{13pt}\selectfont
				Copyrightvermerk:\\
				Dieses Werk einschließlich seiner Teile ist \textbf{urheberrechtlich geschützt}. Jede Verwertung außerhalb der engen Grenzen des Urheberrechtgesetzes ist ohne Zustimmung des Autors unzulässig und strafbar. Das gilt insbesondere für Vervielfältigungen, Übersetzungen, Mikroverfilmungen sowie die Einspeicherung und Verarbeitung in elektronischen Systemen.
			}
		\end{flushleft}
		\begin{flushright}
			{\fontsize{11pt}{13pt}\selectfont \copyright{} 2025 }
		\end{flushright}
	\end{titlepage}
	
\clearpage
\markright{Inhaltsverzeichnis}
\tableofcontents
\newpage
\cleardoublepage
\pagenumbering{arabic}	
\setcounter{page}{1}
% Kapitel Einleitung

\section{Einleitung}
\label{cha: Einleitung}
\subsection{Motivation}
In Zeiten von \textbf{Industrie 4.0} hat die Bedeutung des Automatisierungsgrades in den vergangenen Jahren erheblich zugenommen.\\
Besonders im Hinblick auf Effizienz, Sicherheit und die Effektivität automatisierter Abläufe spielt die Automatisierung eine zentrale Rolle.\\
Der Begriff \textbf{Industrie 4.0} beschreibt den aktuellen Wandel in der industriellen Produktion hin zu intelligent vernetzten Systemen. Dabei spielt die Automatisierung eine entscheidende Rolle. Die Kommunikation der einzelnen Maschinen, Produktionsanlagen und Lager erfolgt hierbei über \textbf{industrielle Netzwerke}, was zu einem echtzeitfähigen Verhalten beitragen kann.\\
Durch die zunehmende Digitalisierung und Vernetzung industrieller Prozesse entstehen stetig neue Anforderungen an die Steuerungs- und Regelungstechnik. Diese Entwicklungen führen zu immer komplexeren, aber auch effizienteren Produktions- und Logistiksystemen.\\\\
Gerade im Bereich der Lagerlogistik ist es entscheidend, sämtliche Sicherheitsmaßnahmen einzuhalten, um einen reibungslosen und sicheren Ablauf zu gewährleisten. Gleichzeitig wird eine präzise und fehlerarme Handhabung von Materialien ermöglicht, was wesentlich zur Prozesssicherheit und Produktqualität beiträgt. Lagerprozesse sind ein wesentlicher Bestandteil des Materialflusses innerhalb der Produktion. Fehler in diesem Bereich können zu Stillständen und hohen Kosten führen. Um solche Probleme zu vermeiden, ist ein reibungsloser Ablauf aller Prozessschritte erforderlich.
\subsection{Bedeutung der Lagerlogistik}
Im Bereich der Automatisierungstechnik existieren zahlreiche Konzepte, um Produkte und Ersatzteile effizient zu lagern und zu verwalten.\\
Eine klassische Variante stellt die manuelle Lagerung in großen Lagerräumen dar. Dabei sind jedoch umfangreiche Such- und Zuordnungsprozesse erforderlich, um die gewünschten Komponenten zu identifizieren und zu entnehmen.\\
Eine wesentlich effektivere Lösung bietet die automatisierte Lagertechnik, insbesondere das sogenannte \enquote{Hochregallager}.\\
Ein Hochregallager ermöglicht die automatisierte Ein- und Auslagerung von Materialien und Ersatzteilen. Solche Systeme ermöglichen es durch den Einsatz moderner Sensorik, wie etwa der Materialerkennung über \textbf{RFID} oder der automatischen Auftragsausführung mittels QR-Code-Erkennung, Materialien präzise zu lagern und zu verwalten. Des Weiteren können die Prozesse über geeignete Softwarelösungen überwacht und gesteuert werden.\\
Dadurch können Prozesszeiten verkürzt, Fehlerquoten reduziert und die Übersichtlichkeit innerhalb des Lagers deutlich verbessert werden.\\\\
Grundsätzlich lässt sich die Handhabung in zwei Varianten unterteilen:
\begin{itemize}
	\item Automatische Ein- und Auslagerung
	\item Manuelle Ein- und Auslagerung
\end{itemize}
Das Hochregallager ist in der Regel in einem separaten Bereich installiert, in dem autorisierte Mitarbeiter mithilfe eines Auftrags- oder Identifikationssystems auf die eingelagerten Komponenten zugreifen können.\\
Die Automatisierung solcher Systeme ermöglicht eine zuverlässige, reproduzierbare und sichere Abwicklung logistischer Prozesse. Hierdurch werden menschliche Fehler weitestgehend reduziert und somit die Sicherheit erhöht.
\subsection{Das DHBW-Hochregallager-Modell}
Zur Simulation und Demonstration der grundlegenden Lagerfunktionen wurde an der Dualen Hochschule Baden-Württemberg (DHBW) eine Miniaturversion eines Hochregallagers entwickelt, die die wesentlichen Prozesse der Ein- und Auslagerung abbildet.\\
Das Modell eignet sich hervorragend als Lern- und Übungsplattform für Studierende, um praxisnah Kenntnisse in der industriellen Automatisierung zu erwerben. Ziel des Systems ist es, die Zusammenarbeit der einzelnen Sensoren und Aktoren zu simulieren und dadurch den Studierenden einen praxisnahen Einblick in die industrielle Automatisierung zu bieten.\\\\
Das System wurde in den vergangenen Jahren bereits mehrfach für Studien- und Projektarbeiten im Studiengang \enquote{Elektrotechnik – Automation} eingesetzt und kontinuierlich weiterentwickelt. Diese Arbeiten ermöglichten es den Studierenden, ein vertieftes Verständnis für Steuerungs- und Automatisierungssysteme zu erlangen und praxisorientierte Erfahrungen zu sammeln.
\subsection{Zielsetzung der Studienarbeit}
Trotz der bisherigen Entwicklungen besteht weiterhin Verbesserungspotenzial, insbesondere im Bereich der Sicherheit und der Programmstruktur.\\
Die vorliegende Studienarbeit befasst sich daher mit der Überarbeitung und Einführung neuer Sicherheitsmaßnahmen. Darüber hinaus soll eine auftragsbasierte Materialein- und -auslagerung implementiert werden. Weitere Schwerpunkte bilden die vollständige Implementierung der \textbf{RFID}-basierten Materialerkennung sowie die Umstellung des bestehenden SPS-Programms auf die Programmiersprache \textbf{Structured Text (ST)}.\\\\
Aufgrund des hohen Umfangs und der technischen Komplexität des Projekts erfolgt die vollständige Umsetzung in den \textbf{Theoriephasen 5 und 6}.\\
Ziel der aktuellen Bearbeitung (Theoriephase 5) ist die Einführung neuer Sicherheitskonzepte sowie die Überarbeitung der RFID-Materialerkennung und des SPS-Programms.\\\\
Abschließend soll die durchgeführte Arbeit die Grundlage für die vollständige Realisierung eines funktionsfähigen, sicheren und didaktisch wertvollen Hochregallagermodells bilden, das zukünftigen Studierenden als praxisnahes Lehr- und Forschungsobjekt dient.\\
Die in den folgenden Kapiteln erläuterten Schritte zur Planung und Umsetzung sollen verdeutlichen, wie die Studienarbeit (Teil 1) zur Weiterentwicklung der Funktionalität und der Sicherheit des Systems beiträgt. Nun erfolgt eine Erläuterung der Problemstellung und des aktuellen Standes der Technik.
%Kapitel Grundlagen 

\section{Grundlagen}
\label{cha: Grundlagen}
\subsection{Problemstellung}
Durch zahlreiche vorangegangene Studienarbeiten wurden bereits viele grundlegende Funktionen des Hochregallagers implementiert.\\
Hierzu gehören unter anderem:
\begin{itemize}
	\item automatisiertes Einlagern,
	\item manuelles Auslagern,
	\item sowie die \textbf{RFID-Materialerkennung}.
\end{itemize}
Diese Funktionen bleiben im Rahmen dieser Arbeit bestehen und sind nicht Hauptbestandteil der weiteren Bearbeitung.\\
Der Fokus der vorliegenden Studienarbeit liegt vielmehr auf der Integration zusätzlicher \textbf{Sicherheitskomponenten}, der Fertigstellung der \textbf{RFID-Materialerkennung} sowie der \textbf{Überarbeitung bzw. Umstellung des SPS-Programms} auf die Programmiersprache \textbf{Structured Text (ST)}.\\\\
Für den sicheren Betrieb eines Hochregallagers sind verschiedene Schutzmechanismen erforderlich, um Gefahren beim Eingriff in den laufenden Prozess zu vermeiden. Einige dieser Maßnahmen wurden bereits in früheren Arbeiten umgesetzt.\\
Ein wesentliches, bislang fehlendes Sicherheitselement ist jedoch eine Lichtschranke, die sich unmittelbar vor dem Lagerbereich befindet. Diese soll im Rahmen dieser Arbeit in das Steuerungsprogramm integriert werden, um potenzielle Gefährdungen bei manuellen Eingriffen während des Betriebs zu verhindern.\\\\
Darüber hinaus ist zur Gewährleistung eines fehlerfreien Ablaufs während der Materialerkennung die Funktionalität der RFID-Prüfung zu analysieren, zu optimieren und gegebenenfalls zu vervollständigen.\\
Ein weiterer wichtiger Aspekt betrifft die verwendete Programmiersprache des SPS-Systems. In den bisherigen Studienarbeiten wurde überwiegend die grafische Programmiersprache \textbf{FBS (Funktionsbausteinsprache)} eingesetzt. Diese bietet den Vorteil einer einfachen, blockorientierten Programmstruktur, stößt jedoch bei komplexeren Anwendungen schnell an ihre Grenzen.\\
Aus Gründen der Übersichtlichkeit, Wartbarkeit und Fehlersuche soll das bestehende Programm daher in die textbasierte Programmiersprache \textbf{ST (Structured Text)} übertragen werden.\\\\
Abschließend lässt sich festhalten, dass die wesentlichen Funktionen des bestehenden Hochregallagers bereits implementiert sind. 
Trotzdem bestehen noch Optimierungspotenziale im Hinblick auf Sicherheit, Struktur und Programmiermethodik. 
Um die geplanten Maßnahmen gezielt umsetzen zu können, ist es erforderlich, den aktuellen technischen Aufbau und die bestehende Programmstruktur zunächst detailliert zu analysieren. 
Im folgenden Abschnitt \enquote{Stand der Technik} werden daher die vorhandenen Komponenten und deren Funktion genauer betrachtet, um eine fundierte Grundlage für die anschließende Überarbeitung zu schaffen.
\subsection{Stand der Technik}





\newpage
\printbibliography
\end{document}